% Ejercicios
\subsection{Ejercicio 1}\label{subsec:ejer1}
Para tener más imágenes con las que trabajar y probar el programa he buscado imágenes libres en la página \textit{Pixabay}: \url{https://pixabay.com/es/photos/search/fotos%20de%20personas/}. A continuación presento las imágenes con las que voy a trabajar:
\begin{itemize}
    \item Figura \ref{fig:imagen1} - Dos personas en un parque. La persona de la izquierda lleva una gorra, tiene barba y está mirando de frente. La persona de la derecha tiene el pelo corto y casi está mirando de frente a la cámara.
    \item Figura \ref{fig:imagen2} - Aparecen ocho personas, todas mirando a cámara y sonriendo.
\end{itemize}

\begin{figure}[ht]
    \centering
    \includegraphics[width=0.8\linewidth]{images/ejer1/imagen1.jpg}
    \caption{Imagen 1}
    \label{fig:imagen1}
\end{figure}

\begin{figure}[ht]
    \centering
    \includegraphics[width=0.8\linewidth]{images/ejer1/imagen2.jpg}
    \caption{Imagen 2}
    \label{fig:imagen2}
\end{figure}

Para esta tarea se ha empleado el clasificador frontal avanzado \texttt{haarcascade\_frontalface\_alt.xml}. También importo la librería \texttt{sys} para poder incluir la imagen con la que se trabaja como parámetro en la orden de ejecución. Esto facilita probar con varias imágenes sin tener que tocar el fichero de código fuente. Además, uso la librería \texttt{os} para guardar las imágenes resultantes (con caras recuadradas) en la carpeta de trabajo.

\begin{lstlisting}[style=pythonstyle, caption={Librerías para el Ejercicio 1}, label={lst:lib1}]
import cv2
import sys
import os
\end{lstlisting}

Como la imagen con la que se va a trabajar se pasa por parámetro hay que comprobar que se haya pasado, que la ruta exista y que no esté corrupta:
\begin{lstlisting}[style=pythonstyle, caption={Comprobación de parametro}, label={lst:param1}]
# Validación del número de argumentos en la línea de comandos
if len(sys.argv) < 2:
    print("[ERROR] Uso correcto: python ejercicio1.py <nombre_imagen.jpg>")
    sys.exit(1)

# Captura del parámetro del archivo de imagen
nombre_imagen = sys.argv[1]

#Cargamos nuestro modelo classificador de Haar:
modelo = cv2.CascadeClassifier('haarcascade_frontalface_alt.xml')
  
#Cargamos la imagen
img = cv2.imread(nombre_imagen)

# Control de excepciones ante rutas inexistentes o corruptas
if img is None:
    print(f"[ERROR] No se pudo cargar la imagen '{nombre_imagen}'")
    sys.exit(1)    
\end{lstlisting}

A continuación, el programa sigue como el del ejemplo, haciendo el procesamiento de escala de grises, buscando rostros y dibujando los rectángulos sobre la imagen original. Una vez hecho eso, se guarda la imagen en disco separando el nombre de la imagen de la extensión:
\begin{lstlisting}[style=pythonstyle, caption={Guardar imagen}, label={lst:guardarimg1}]
# Separamos el nombre del archivo de su extensión (ej: 'foto' y '.jpg')
nombre_base, extension = os.path.splitext(nombre_imagen)

# Creamos el nuevo nombre para el archivo de salida
nombre_salida = f"resultado_{nombre_base}{extension}"

# Guardamos la imagen modificada en la carpeta de trabajo
cv2.imwrite(nombre_salida, img)
print(f"[INFO] Imagen resultante guardada con éxito como: {nombre_salida}")
\end{lstlisting}

El resultado usando \texttt{scaleFactor=1.35, minNeighbors=6} en la Figura \ref{fig:imagen1} es la Figura \ref{fig:1res_img1}. El motivo de usar \texttt{scaleFactor=1.35} es porque con el valor $1.5$ no se reconoce ningún rostro. El rostro del hombre no lo llega a reconocer aunque se disminuya el valor de \texttt{scaleFactor}.\\

\begin{figure}[ht]
    \centering
    \includegraphics[width=0.8\linewidth]{images/ejer1/resultado_imagen1.jpg}
    \caption{Resultado de \textit{imagen1} con \texttt{scaleFactor=1.35, minNeighbors=6}}
    \label{fig:1res_img1}
\end{figure}

El resultado usando \texttt{scaleFactor=1.5, minNeighbors=6} en la Figura \ref{fig:imagen2} es la Figura \ref{fig:1res_img2}. Se reconocen 5 rostros, de 8 que hay en la imagen.\\

\begin{figure}[ht]
    \centering
    \includegraphics[width=0.8\linewidth]{images/ejer1/resultado_imagen2.jpg}
    \caption{Resultado de \textit{imagen2} con \texttt{scaleFactor=1.5, minNeighbors=6}}
    \label{fig:1res_img2}
\end{figure}

\begin{figure}[ht]
    \centering
    \includegraphics[width=0.8\linewidth]{images/ejer1/resultado_imagen2_2.jpg}
    \caption{Resultado de \textit{imagen2} con \texttt{scaleFactor=1.35, minNeighbors=6}}
    \label{fig:1res_img2_2}
\end{figure}
Si disminuimos el valor de \texttt{scaleFactor} a $1.35$ el programa reconocerá todas las caras, como se muestra en la Figura \ref{fig:1res_img2_2}.\\
